\section*{G. A published learned repair tested for cooperative admission}
We additionally test the preferred learned Arbiter solution in
Appendix~C.4.1 of~\cite{degradation}.
The same declared reconstruction from the eleven-candidate cohort restores the omitted response-pattern declaration: 886 bytes, no imports. After inspecting prior results, we selected the smallest candidate with empty original $V,K$ among those previously positive under residual initialization. This outcome-informed case is separate from
the fixed 14-input public-corpus comparison. Cooperative admission adds to the degradation study's ordinary-realizability objective, which we do not dispute or historically reproduce.

\paragraph{Source obstruction.}
The environment fields are $a,r_1,r_2$, the public system fields are
$g_1,g_2$, and two system-owned response monitors are retained. The
original environment safety is
\[
 E(s)=a\lor(\neg g_1\land\neg g_2\land\neg r_1\land\neg r_2),
 \qquad \rho_e(s,e')=E(s),\quad\theta_e=\mathit{true}.
\]
Because this current-only invariant refers to system variables, the
audited StateInvTranslator leaves it as current safety; it neither
strengthens the environment initializer nor globally primes the formula.
Ordinary system invariants contribute current initialization and next-state
safety. The response monitors update from current requests and grants.

If $E(s)$ is false, no legal environment input exists. Otherwise every
next input is legal, including $(a',r_1',r_2')=(0,1,0)$. If no safe
response exists, strict winning fails. Every safe response instead reaches
$t$ with $E(t)=\mathit{false}$, regardless of system values, auxiliary
states, or private memory. That legal prefix has no infinite legal
continuation, hence no assumption-fair continuation. Thus no residual
state admits a strict cooperative strategy. The nonempty original
initializer then excludes initial cooperative realizability; this does
not assert failure for artificial vacuous initializers.

\paragraph{Complete output and independent check.}
The configured released job returns a complete 112-state, 1,176-edge
native graph with eight initial states covering all initial environment choices. The reachable helper-false
view contains 56 states and 280 edges, preserving native IDs, rank memory,
both monitors, and initial choices.
Its original safety and input coverage pass; strict liveness holds,
but 21 reachable states have no assumption-fair continuation.
An independent implementation directly evaluates the saved decision DAGs
against source-coded equations on all 128 original states and 16,384
state pairs. It checks the native printer, ordered justices, 1,792 native
input candidates and 448 original-view input candidates. Bitset transitive
closure independently establishes the graph properties. It imports no
production BDD or fairness checker; its result was sealed before comparison
with that checker, after which the state, edge, justice and SCC sets agreed.
Both checks share compiler exports; this source-equation check is specific to the contract, not a general compiler proof.

\paragraph{Two attempts and limits.}
The first native attempt stopped at an 8-MiB stdout cap and remains UNKNOWN.
A separately recorded followup kept the source, runtime, classes, synthesis
settings, extractor, printer and checker unchanged. Native stream caps
rose from 8 to 64 MiB, the wall bound fell from 60 to 10 seconds, and
the output postcheck rose from 64 to 128 MiB; checker limits stayed fixed.
It completed with all 11,927,997 stdout bytes retained, normal teardown
and child reaping: one complete result from two attempts on one candidate.
No transformed strict-winning region was calculated. Neither native
acceptance nor the complete-literal $C=D$ result establishes preservation
in this released arena. No repair learning, adaptation loop or deployment
was executed; other cohort counts remain unchanged.

\section*{H. The fixed eleven repair candidates and selected Lift behavior}
\label{supp:repair-eleven}
The completed comparison covers all eleven replacements in
Appendix~C.3--C.4 of~\cite{degradation}: five preferred learned candidates
and six trivial alternatives from five examples. Two Minepump alternatives
use the unchanged listings. The other nine use the declared executable
reconstructions appending the paper's omitted response-pattern declaration;
Traffic Single P also lowers its direct eventuality to that pattern.
Lift's one-floor assumptions and movement guarantees are unchanged from
the published replacement listings. Its appended pattern includes auxiliary
initialization, deterministic updates, and justice~\cite[Appendix C.2.1, C.3.3, and C.4.3]{degradation}.
These are correlated published candidates, not eleven independent
applications. The constructed controls in Section~I are excluded.
The extraction intervention and modeled episode below are later
followups on existing inputs; they do not alter this cohort's
released-route outcomes or denominator.

\paragraph{Traffic Single's direct-eventuality lowering.}
The learned Traffic Single guarantee $\G(s\Rightarrow\F p)$, with
$s=\mathit{car}$ and $p=\mathit{green}$, uses the paper's response pattern.
Writing $q=0$ for $S0$ and $q=1$ for $S1$, its deterministic monitor is
$q_0=0$, $q'=(q\lor s)\land\neg p$, with justice $\G\F\neg q$.
Every infinite observable word has exactly one monitor run.
If every trigger receives a current or later response, each pending
obligation clears; either clearances recur infinitely or $q$ eventually
stays zero. Conversely, an unanswered trigger makes $q$ permanently
true from the next step, violating the justice.
Hence the two forms have the same infinite observable trace language
after projecting out $q$. This is the saved reconstruction argument;
compiler correspondence, auxiliary ownership, and Mealy support remain
separate obligations. The guarantee retains SYS ownership.
The equivalence starts at the original initializer; continuation checks
retain the reached monitor value.

\paragraph{Initialization, correspondence, and output checks.}
For contract-feasibility checks, every case uses its original initializer,
including all logical auxiliaries.
For the independently certified original cooperative region $K$, we check
\[
 I_M(K)=\forall e:\theta_e(e)\Rightarrow
              \exists y:\theta_s(e,y)\land K(e,y),
\]
where $y$ includes the logical system monitors. All eleven have nonempty
legal initial environment domains; eight pass this predicate. The saved
region authority uses the original nonblocking arena and the normalized
round/alternating correspondence of Section~A. Neither a residual trace
entry nor $K\ne\varnothing$ substitutes for initialization.

\paragraph{Minepump P: strict initialization and the native refusal.}
The saved original strict region contains 28 of the 32 typed states.
Its three-node Boolean representation is
\[
 W=\neg\mathit{highwater}\lor\neg\mathit{methane}\lor\mathit{pump}.
\]
The original environment initializer fixes highwater and methane to
false; the system initializer fixes pump to false, the response monitor
to $S0$, and the previous-pump auxiliary to false. Substitution into the
saved $W$ gives $I_M(W)=\mathit{true}$, while the existing nonblocking
and cooperative certificates give $V=K=\varnothing$ and $I_M(K)=\mathit{false}$.
This is a static reading of the saved region and initial predicates,
not a new solver run. The historical residual trace entry has
highwater and methane true and pump false, so it lies outside $W$;
it is not an original initial state. The configured route's native
False cannot be transferred as a claim of original strict
unrealizability. Cooperative rejection instead follows from the empty
$V$; no complete current released strict-arena explanation of that raw
refusal is asserted.

The released route and settings are specified in Sections~L.3 and T;
Section~Y.1 distinguishes that selected route from the strict diagnostic
queries. All eleven
adopted candidate jobs retain the original initializers; earlier residual
trace entries describe separate studies. The nine complete correspondence-
gated views cover all original ENV initial choices. Lift P's direct graph
audit and Minepump P's no-controller outcome are distinguished below. Complete
native outputs are preserved in the study records. The Lift static addendum
omits Lift P's historical 220,856,594-byte stdout, so that stream cannot
be independently rehashed from the addendum. The correspondence-gated audit requires typed equality
of initial, safety, auxiliary-update and ordered-justice predicates under
identity or the helper-false mapping. The environment initializer is also
checked over the full native typed domain. For the nine newly executed
candidates, separate implementations check admitted graphs for original initial inputs,
every legal next input, the selected safe response
and private-memory successor, then strictness and cooperation. Saved
comparisons require equal state/start/edge sets, ordered justice labels,
SCC evidence and retained raw nodes, including ranks and auxiliaries.
A failed correspondence or incomplete output remains UNKNOWN at this gate.

Independent hand-derived source equations are also compared with saved
initialization, safety, and ordered justice predicates over the full
typed domains, retaining logical auxiliaries. For Arbiter P and Traffic S P,
the prior audits use the helper-false view of the released model;
a separate nine-candidate batch checks helper-free originals for every
remaining candidate. These comparisons use the documented translation
conventions. Source-to-original equality does not resolve Lift P's
distinct released-model mismatch below.

\begin{table}[htbp]
\centering\small
\caption{Correspondence-gated audit of the fixed eleven candidates. P denotes preferred/learned, T trivial,
and S/U Traffic Single/Updated. Native False is the solver's negative
verdict; dashes mean no controller. Question marks mean original
correspondence failed. $^{b}$ denotes a reused prior baseline and
$^{f}$ a separate Lift resource followup. This table retains the released-route baseline; the later extraction intervention and the post-hoc direct Lift P audit below are separate.}
\label{tab:supp-repair-eleven}
\begin{tabular}{@{}lccrcc@{}}
\toprule
Candidate & $I_M(K)$ & Native & View states/edges & Strict & Coop.\\
\midrule
Arbiter P$^{b}$ & F & Graph & 56/280 & Pass & Fail\\
Arbiter T & T & Graph & 56/448 & Pass & Pass\\
Lift P$^{f}$ & F & Graph & ? & ? & ?\\
Lift T$^{f}$ & T & Graph & 154/1170 & Pass & Fail\\
Minepump P & F & False & --- & --- & ---\\
Minepump T1 & T & Graph & 12/40 & Pass & Pass\\
Minepump T2 & T & Graph & 4/16 & Pass & Pass\\
Traffic S P$^{b}$ & T & Graph & 16/80 & Pass & Pass\\
Traffic S T & T & Graph & 12/80 & Pass & Pass\\
Traffic U P & T & Graph & 68/323 & Pass & Pass\\
Traffic U T & T & Graph & 56/416 & Pass & Pass\\
\bottomrule
\end{tabular}
\end{table}

\paragraph{Attempt history and denominator.}
Arbiter P reuses the complete followup in Section~G; Traffic S P reuses
its earlier 52-state/496-edge native baseline. Exact certificates and
source-byte mappings are retained. Each remaining candidate received
one declared native attempt. Six returned graphs passing both checks,
Minepump P returned native false, and both Lift attempts stopped at the
64-MiB stdout cap, not the ten-second wall bound. Those two attempts
remain resource UNKNOWN.

One separately authorized additional attempt per Lift candidate kept the
sources, runtime, synthesis settings, extractor and audits fixed. Stream
caps rose to 256 MiB and the per-case output threshold from 128 to
512 MiB; the ten-second native wall and 1-GiB RSS limits were unchanged.
Both completed with full stdout: 220,856,594 bytes for Lift P and
219,540,020 for Lift T. The former still fails original correspondence;
the latter has a certified strict-only view. At this gate, ten candidates return
native graphs, nine have complete original-view certificates, seven pass
both properties, two pass strictness only, and one returns native false.
Followups do not add candidates or rewrite prior attempts.

\paragraph{A separate original-initializer synthesis route.}
\label{supp:normalized-eleven}
After recording the historical outcomes above, we froze the same eleven
source files, their order, and the settings for one additional route.
It applies original-$V$ restriction and justice normalization
$G\mapsto G\land A$, then the existing normalized four-fixed-point
solver and unchanged generic graph constructor. It adds no helper and
does not apply the final-choice preference used in the Lift intervention.
The source initializers remain unchanged: every legal initial ENV choice
must have a SYS/AUX response satisfying $\theta_s\land K$, and extracted
controllers initialize their private rank to zero. Saved replay or
residual entries are not substituted for these initial conditions.
All eleven source ENV and joint initializers are nonempty.

Each source received one serial native process attempt, with concurrency
one and zero retries. Eleven attempts led to eight solver/extractor calls;
Arbiter P, Lift P, and Minepump P stopped before those calls on independently
confirmed empty $V$. All eleven original arenas, including initializers,
safety, typed auxiliaries, and ordered justices, match their source-bound
saved comparators. The full $V$ and $K$ Boolean functions also agree,
including all AUX valuations. The comparisons use exact exported predicates
and shared parser-translated BDDs; they are not an independent compilation.

\begin{table}[htbp]
\centering\small
\caption{Separate normalized-route results from the original initializer,
in the fixed execution order. Init. counts graph initial nodes. Coop. is
legal-prefix cooperation. Empty $V$ is an independently checked refusal
before solver and extraction; its controller flags are null, shown as
dashes. This table does not replace Table~\ref{tab:supp-repair-eleven}.}
\label{tab:supp-normalized-eleven}
\begin{tabular}{@{}lclrrcc@{}}
\toprule
Candidate & $I_M(K)$ & Outcome & States/edges & Init. & Strict & Coop.\\
\midrule
Arbiter P & F & Empty $V$ & --- & --- & --- & ---\\
Arbiter T & T & Graph & 56/448 & 8 & Pass & Pass\\
Lift P & F & Empty $V$ & --- & --- & --- & ---\\
Lift T & T & Graph & 154/1176 & 2 & Pass & Pass\\
Minepump P & F & Empty $V$ & --- & --- & --- & ---\\
Minepump T1 & T & Graph & 12/40 & 1 & Pass & Pass\\
Minepump T2 & T & Graph & 4/16 & 1 & Pass & Pass\\
Traffic S P & T & Graph & 28/152 & 8 & Pass & Pass\\
Traffic S T & T & Graph & 20/120 & 8 & Pass & Pass\\
Traffic U P & T & Graph & 116/546 & 8 & Pass & Pass\\
Traffic U T & T & Graph & 84/640 & 8 & Pass & Pass\\
\bottomrule
\end{tabular}
\end{table}

All eight complete graphs pass the original-contract checks for initial
coverage, legal-input coverage, safe responses, private-memory successors,
strictness, and legal-prefix cooperation. The three refusals have false
$I_M(K)$ and no controller verdict. This run has
no UNKNOWN or unattempted row. Its graph certificates come from one
original-contract checking implementation; the two-judge Lift experiment
and the separate fourteen-input comparison retain their own scopes.

\paragraph{Interpretation and limits.}
The baseline and normalized jobs use the same original initializers.
The new eight accepted graphs and historical seven passing outputs come
from separate whole-route executions, not a matched extraction or
success-rate comparison. Historical refusals, resource UNKNOWNs, and
Lift P's released-model correspondence failure remain unchanged.
For Lift T, the new certificate records zero copied source ENV justices
and a 154-state, 1,176-edge graph with two initial nodes. Its equal size
does not establish identity with the intervention graph; no graph-equivalence
comparison was performed. The successful route supplies an existing
alternative without the new preference, while the paired same-game,
same-memory intervention still shows one sufficient local change.
This whole-route comparison does not isolate the effects of $V$ restriction,
normalization, helper removal, or extraction, and establishes
neither preference necessity nor a general synthesis failure.
It does not test Friendly's specialized strategy construction or add an
eleven-case Slugs run. No new route output was installed or executed in
the maintenance-to-service episode. The fixed correlated sources and
single attempts establish no timing superiority or broader generality.


\paragraph{Source clauses shared by Lift T and Lift P.}
The Boolean environment inputs are $b_1,b_2,b_3,c$ and the system
outputs are $f_1,f_2,f_3$. Write $m_i$ for the SYS-owned logical
response monitor of button $i$, encoding $S0=0$ and $S1=1$.
Let $\operatorname{One}(f)$ mean exactly one true floor. The common
environment safety relation is
\[
E_L=\operatorname{One}(f)\land
\bigwedge_{i=1}^{3}\left[
(b_i\land f_i\Rightarrow\neg b_i')\land
(b_i\land\neg f_i\land\neg c\Rightarrow b_i')\right].
\]
One-floor consistency is thus a condition on the \emph{current}
outputs in environment safety. The movement guarantees are
$\bigwedge_i(f_i\Rightarrow D_i)$, where
\[
\begin{aligned}
D_1={}&(f_1'\land\neg b_2\land\neg b_3)
 \lor(f_2'\land b_2)\lor(f_2'\land b_3)\lor(f_1'\land c),\\
D_2={}&(f_1'\land b_1)
 \lor(f_2'\land\neg b_1\land\neg b_3)
 \lor(f_3'\land b_3)\lor(f_2'\land c),\\
D_3={}&(f_3'\land\neg b_1\land\neg b_2)
 \lor(f_2'\land b_1)\lor(f_2'\land b_2)\lor(f_3'\land c).
\end{aligned}
\]
Every button and cancellation guard here uses the current value.
System safety conjoins these movement implications with
\[
\bigwedge_{i=1}^{3}
\left[m_i'\leftrightarrow(m_i\lor b_i)\land\neg(f_i\lor c)\right].
\]
In particular, it does not add $\operatorname{One}(f')$: a movement
disjunct can permit extra true next floors. Initialization sets
$b=000$, $f=100$, and $m=000$, leaving $c$ unrestricted.
The environment justice list is empty; the three system justices are
$\mathbf{G}\mathbf{F}\neg m_i$. Lift T has $\rho_e=E_L$.
These are the clauses of \path{lift__trivial_1.spectra}, lines 3--81,
shared with \path{lift__preferred.spectra} except for the latter's
two added assumptions stated below. Supplement K's Lift T entry links
the source, policy, positive graph, complete extension, and selected graph.
The policy's private priority bit is separate from all three logical monitors.



\paragraph{Lift T: feasibility and a failing selected graph.}
The typed original carrier has 1,024 states, including three response
monitors $m_i\in\{S0,S1\}$. The saved regions have $|V|=384$ and
\mbox{$|K|=316$}. Both legal initial choices, distinguished by cancellation
input $c$, have the unique initialized system/monitor response in $K$.
The native graph has 794 states and 13,874 edges; its original view has
154 states, 1,170 edges and two initial nodes. Both checkers establish
strictness and reject cooperation.

Table~\ref{tab:lift-t-prefix} records the original-legal prefix
$0\to10\to339\to438\to704\to680$.
It first leaves both $K$ and $V$ on its last edge.
All five prefix edges satisfy original safety and the exact current-state
monitor updates. The final move is allowed by the next-$f_3$ disjunct
of the current floor-2 move constraint. Its resulting two true floors
violate the environment's current-state one-floor assumptions, making
every next original environment input illegal. Nodes 680 and 708 are
the only two original-view nodes outside both $K$ and $V$; each has
zero original-view successors and 16 native successors, all helper-true.
The original environment justice list is empty, but a finite dead end
still fails the infinite-extension requirement.

\begin{table}[ht]
\centering\small
\caption{Recorded original-legal Lift T prefix, starting at native
initial node 0. Bit strings use index order 1,2,3, with $S0=0,S1=1$;
the helper is false throughout. The recorded native system-justice
index $r$ names $\mathbf{G}\mathbf{F}\neg m_{r+1}$,
not a decreasing proof rank.}
\label{tab:lift-t-prefix}
\begin{tabular}{@{}rccccr@{}}
\toprule
Node & $b_1b_2b_3$ & $c$ & $f_1f_2f_3$ & $m_1m_2m_3$ & $r$\\
\midrule
0   & $\mathtt{000}$ & 0 & $\mathtt{100}$ & $\mathtt{000}$ & 0\\
10  & $\mathtt{010}$ & 1 & $\mathtt{100}$ & $\mathtt{000}$ & 0\\
339 & $\mathtt{001}$ & 1 & $\mathtt{010}$ & $\mathtt{000}$ & 0\\
438 & $\mathtt{100}$ & 0 & $\mathtt{001}$ & $\mathtt{000}$ & 1\\
704 & $\mathtt{101}$ & 0 & $\mathtt{010}$ & $\mathtt{100}$ & 2\\
680 & $\mathtt{101}$ & 0 & $\mathtt{011}$ & $\mathtt{101}$ & 0\\
\bottomrule
\end{tabular}
\end{table}

For the same observed next environment tuple and the same required next
monitor tuple $(S1,S0,S1)$, changing only next $f_2$ to false gives
$(f_1',f_2',f_3')=(0,0,1)$. This satisfies source safety and lies in
saved $K\cap V$. It is not a selected successor of node 704, and we
assign it no replacement native rank or private memory. This one-step
alternative is not a repaired controller or a composed strategy.
It is also distinct from the later extraction experiment: at first
creation of a bad node, the final native candidate retains floor
$100$, whereas the nearest-$Y$ progress filter in rank selection has
already removed $001$. On the displayed
$704\to680$ edge, reuse of an existing node has already removed every
continuable candidate. The diagnosis below distinguishes these stages.

\paragraph{A separate positive source policy.}
A fixed, independently checked policy directly witnesses original initial
feasibility. It always sets exactly one floor true and uses one private
priority bit $d\in\{1,3\}$, initially 3. From floor 1 or 3 it moves to
floor 2 exactly when another floor is requested. At floor 2 it chooses
a requested endpoint, stays if neither is requested, and uses $d$ if
both are requested; a departure toward $j$ sets $d'=4-j$. Decisions use
current buttons. Encoding $S1$ as true, every logical monitor updates as
$m_i'=(m_i\lor b_i)\land\neg(f_i\lor c)$.

Every move satisfies an explicit source disjunct. The one-floor invariant
keeps legal environment successors available. Take a state at time $t$
with $b_i$ true. If $f_i$ or $c$ already holds, the request is answered.
Otherwise its button persists until arrival or cancellation. An endpoint
request can be delayed by at most one trip to the opposite endpoint
before priority reverses, so the requested floor is reached within four
transitions from $t$ unless cancellation occurs first. Arrival or
cancellation makes its monitor $S0$ in the next state, because its
update reads the current $f_i$ and $c$. All three monitors therefore visit $S0$ infinitely often. Exhaustive
construction from both original initial states (both values of $c$, with
all three source monitors initially $S0$) produces 128 nodes and
1,020 edges, preserving every legal input, all ten original fields and
the private bit. Direct source and independent finite-graph checks certify
strictness and cooperation. This separate source-coded existence witness
uses no native graph, rank, saved $K/V$ restriction or repair as a
construction input; it makes no claim from arbitrary orphan monitor
states or every state of $K$.

\paragraph{Extension to the complete transformed model.}
For an explicit cooperative suffix after any legal policy prefix, keep
$c'=1$ and choose $b_i'=b_i\land\neg f_i\land\neg c$ at every step.
The two ENV button implications are satisfied, the policy preserves
exactly one floor, and current cancellation clears all monitors within
two transitions. Continuing these choices gives an infinite ENV-safe
suffix; the source has no ENV justices. This is a semantic witness,
not a claim that cancellation was used in the recorded workloads.

The same positive policy also extends to a strict strategy for the
complete saved transformed model, with environment helper $h$. Keeping its exact source AUX updates
and private priority, we checked all 49,152 prescribed transitions over
every one-hot current floor, public input, AUX valuation, priority and
current/next helper value. Every response satisfies both original system
safety and complete native system safety. This includes all first-flip
and post-flip inputs. On this one-hot carrier, which the policy preserves,
the helper cannot revert on a native-legal edge.
Initialization, domains and ordered justices were checked on all 2,048
native typed states. We do not approximate native system safety by
merely disjoining the helper into the source relation, because AUX
updates remain required.

Starting from either original initial input, a play that keeps the helper
false is an execution of the certified 128-state source controller.
If the helper rises, it remains true, falsifying the environment justice
$\mathbf{G}\mathbf{F}\neg h$, while system safety continues. Thus the
policy is native-strict and its helper-false view is cooperative. This
establishes neither liveness from arbitrary orphan-AUX states nor
cooperation after a helper flip. This source-policy argument alone does
not attribute the selected controller's failure; the following paired
experiment supplies separate evidence about extraction choices.

\subsection*{Lift T: controlled extraction and modeled replacement}
\label{supp:lift-causal}
For this Lift T instance, changing one extraction preference while
holding the solved game and fixed-point memory fixed produces a strict,
cooperative original-contract controller. A subsequent constructed
execution uses the checked controller after an environment-model
weakening. Both are followups designed after inspecting the baseline
failure, not held-out validation or a general synthesis correction.


\paragraph{Baseline and fixed-point agreement.}
A fresh run reproduces all 23 saved semantic artifacts of the released
baseline, including its complete 794-state, 13,874-edge native graph and
154-state, 1,170-edge original-contract view. Both graph auditors again
accept strictness and reject cooperation. A behavior-preserving observer
uses the same model and memory objects; complete graph equality,
semantic equality of the saved model and memory predicates, and
unchanged variable order establish observation parity.

A separate finite-set implementation evaluates the inspected raw
four-fixed-point recurrence on all 2,048 transformed typed states.
It uses the saved transition predicates and their original move
quantifiers, without an \emph{Alive} restriction or justice normalization.
Its region contains all 2,048 states, exactly matching the native region.
The effective 12 $X$ slices, 12 $Y$ slices, three $Z$ components and final
winning set also agree on all 57,344 state memberships.
The recurrence uses Python integer sets; a shared saved-BDD decoder
only reads predicate membership. These checks establish numerical
agreement on this input, not compiler independence, the raw core's
semantic soundness premise, or a general progress-memory theorem.

\paragraph{The transformed strict game permits the dead end.}
Node 680 belongs to the raw region and satisfies the original
response condition $C$. A focused check of its saved residual strategy
also witnesses membership in the actual transformed strict-winning
region $W_R(\cdot,0)$. All 16 legal next inputs flip the helper;
none can keep it false. The residual graph has 321 states and
5,136 edges. At every reached node it covers exactly every legal
input with a unique safe response, retaining all monitors and private
memory. Every successor has the helper true and violates the added
environment justice, so every infinite legal play satisfies the strict
liveness implication. This uses the complete saved relations, rather
than assuming that the syntactic arena matches the complete-literal
theorem. The check reuses the saved-BDD decoder; it performs no new
native generation or region solution.
Thus the game itself permits this non-cooperative behavior. The paired
intervention below shows that the same upstream data also permit an
accepted output; it does not identify extraction as the unique cause.

\paragraph{Where the continuable alternative is lost.}
Here the original environment has no justice, and
$\mathit{Alive}(s)=\exists e':\rho_e(s,e')$ equals the current one-floor
condition. The trace records two first creations of bad nodes and two
later reuses. At first creation, rule~a's raw successors have six floor
patterns, including the live $001$ and $100$ alternatives. After its
nearest-$Y$ choice, the final five-pattern candidate still contains
$100$ but not $001$; \code{satOne} selects the dead $011$.
Later rule~a prefers the existing dead node before the final choice.
Thus the source-level $001$ alternative above and the final-stage $100$
alternative have different roles. The trace contains no dropped events
or predicate references; it does not by itself certify a fair
continuation after any counterfactual choice.

\paragraph{Single-axis intervention and complete-controller checks.}
Immediately before transition \code{satOne}, the experimental constructor
intersects its final candidate with an actual helper-false input-availability
preference. It uses this intersection when nonempty, otherwise retaining
the native candidate. With helper $h$, the preference is
\[
 h(s)\ \lor\
 \exists e'\,[D_e(e')\land
       \rho_e^{\mathrm{native}}(s,e')[h=0,h'=0]] .
\]
Without a helper it uses the supplied native environment relation.
The helper-false relation is not universally identical to the original
environment relation: Section~N gives a static translation counterexample.
Lift's complete original-contract graph checks below establish its
instance-specific result independently of that universal interpretation.
This leaves helper-true candidates available. It changes no source
clause, initializer, solver, stored $X/Y/Z$ set, variable order, rank
rule, or earlier existing-state preference. There is no hardcoded floor
constraint or node ID.

One solver call supplies the same model and memory objects to the
released and experimental extractions. Their initial, transition,
ordered-justice and effective memory predicates remain semantically
equal; the baseline again matches the saved graph. Both complete outputs
are checked by the production and separately implemented graph auditors,
whose mathematical checks are unchanged and whose metadata gates
explicitly identify the nonrelease variant.
The production reader reconstructs saved predicates as BDDs; the other
auditor evaluates the exported decision DAG with a separate reader and
compares it to separately encoded source equations. Each run checks all
$1{,}048{,}576$ original current/next state pairs for correspondence and
$2{,}464$ typed input obligations on its 154-node original view.
The second auditor imports neither the production checker nor its
certificate and uses a separate reachability-based fairness test.
Both still depend on the same compiler exports, typing information,
and recorded translation convention; implementations were mutually
visible. The shared decoder limitations for the fixed-point and
residual-strategy checks above are distinct from this graph-audit pair.
Table~\ref{tab:lift-choice-intervention} reports the original-contract
views. The experimental full native graph has 778 states and 13,624
edges; 726 of its 13,624 final candidates are refined, with no
empty-preference fallback. This is one code-level change affecting many
choices, not four patched witness edges.

\begin{table}[htbp]
\centering\small
\caption{Paired extraction results on Lift T and six prespecified
successful controls. Each cell gives original-view states/edges and
strict/cooperative verdicts; P/F mean pass/fail. Refinements count
candidate sets changed by the preference. These are overlapping
followups, not additional members of the eleven-candidate cohort.}
\label{tab:lift-choice-intervention}
\begin{tabular}{@{}lccc@{}}
\toprule
Input & Released & Experimental & Refinements\\
\midrule
Lift T & 154/1170, P/F & 154/1176, P/P & 726\\
Arbiter T & 56/448, P/P & 56/448, P/P & 0\\
Minepump T1 & 12/40, P/P & 12/40, P/P & 0\\
Minepump T2 & 4/16, P/P & 4/16, P/P & 0\\
Traffic S T & 12/80, P/P & 12/80, P/P & 0\\
Traffic U P & 68/323, P/P & 68/323, P/P & 0\\
Traffic U T & 56/416, P/P & 56/416, P/P & 0\\
\bottomrule
\end{tabular}
\end{table}

The six controls were declared before their new executions and retain
both verdicts with zero refinements and zero fallbacks. They establish
preservation of those successes when the preference is dormant, not
generality when it is active. For Lift, unchanged game, raw region and
memory suffice for this constructor to return a checked cooperative
view after changing concrete extraction choices. One-step availability alone is no
fairness proof; the positive conclusion comes from complete-controller
checks. No sole-defect, missing-$W$ cause, or general synthesis-repair
claim follows.

\paragraph{A constructed environment change and exact handoff.}
The old contract is Lift T with the additional environment restriction
that all current and next buttons are false. A source-coded old
controller stays at floor $100$ and performs the actual response-monitor
updates. Its two reachable cancellation states and four edges pass
strictness and cooperation under this maintenance contract.
Two modeled plants execute four old-controller steps, with next
cancellation values $0,1,1,0$. Removing only the maintenance restriction
then exposes the original Lift T environment. Both plants pause at
the existing state $b=000,c=0,f=100,m=000$.

After this change, a fresh paired native run generates both candidates,
followed by four fresh graph judgments; no cached candidate is installed.
The fixed order is released controller, then experimental controller.
One admission policy accepts the first strict candidate; the other
requires both strictness and cooperation. Each also requires exact
agreement between the current full source valuation and an allowed
initial native node. Installation preserves every source field and
logical monitor and initializes the new controller's private memory to
that matching node. The handoff therefore uses an already
initial-compatible paused state; it is not arbitrary-state migration
or a reset of pending requests.

\paragraph{Executed outcomes.}
Both arms execute the five inputs of the known failing prefix,
in $(b,c)$ notation:
$(010,1),(001,1),(100,0),(101,0),(101,0)$.
The strict-only policy installs the released controller and reaches
floor $011$ after these five post-change steps. Independent enumeration
finds all 16 possible original next inputs illegal.
The strict-and-cooperative policy rejects that candidate, installs the
experimental one, and reaches floor $100$ instead.
It then executes the lexicographically smallest legal input at each
step, with false ordered before true in $(b_1,b_2,b_3,c)$.
Within the declared 32-step continuation bound, four more transitions
give native nodes
$680\to524\to544\to434\to434$.
The final self-loop has $b=000,c=0,f=001,m=000$, meeting all three
system justices; there are no original environment justices.
Native IDs are local to each graph: experimental node 680 is not the
released graph's dead node 680.

A separate no-replacement counterfactual applies $(010,0)$ twice to the
stationary old controller under the weakened contract. The first step
is legal; the second proposed response violates the movement guarantee
and is rejected. Thus keeping this old controller does not solve the
changed model.

The independent episode audit imports neither the episode runner nor
its transition implementation. It checks the 47-event log, source
semantics, exact handoff, installed-controller hashes and actual native
responses, and binds the four fresh full-graph judgments. It performs
24 environment and 24 system source checks, including the rejected
counterfactual step, and 42 native-edge/response checks. All six
supervised processes, comprising the outer execution, one native pair
and four judges, complete with verified cleanup.
The executed loop witnesses one infinite continuation; the separate
complete-graph certificates justify the controller properties.
This is a constructed in-process episode with a known-failure prefix,
not hardware deployment, newly learned repair, arbitrary memory
transfer, or a runtime-cost benchmark.

\paragraph{Bounds, attempts, and provenance.}
Table~\ref{tab:lift-followup-bounds} gives the declared successful-run
limits. Stream bounds apply separately to stdout and stderr.
RSS caps apply to the owned process groups, not a measured aggregate
tree; output bounds are sampled termination thresholds with final
checks, not exact quotas. The episode's outer final \code{wait4}
resource value can include waited descendants on this host.
Native and checker children retain a 16-MiB per-artifact limit; the
outer episode's per-file limit is 1~GiB so it can capture the already
bounded native streams. Each launch requires 1~GiB of free disk space.

\begin{table}[htbp]
\centering\small
\caption{Bounds for the accepted diagnostic and episode runs.
The episode includes serial child processes with their own limits.}
\label{tab:lift-followup-bounds}
\begin{tabular}{@{}lrrrr@{}}
\toprule
Process & Timeout (s) & RSS (GiB) & Stream (MiB) & Output (MiB)\\
\midrule
Fresh baseline & 10 & 1 & 256 & 512\\
Observer / paired extraction & 30 & 1 & 512 & 768\\
Finite-set / memory comparison & 120 & 1 & 1 & 32\\
Production graph judge & 60 & 1 & 8 & 64\\
Independent graph judge & 180 & 1 & 8 & 64\\
Outer modeled episode & 600 & 2 & 8 & 1024\\
\bottomrule
\end{tabular}
\end{table}

\paragraph{Observed durations of the completed episode.}
Table~\ref{tab:lift-episode-observed-times} extracts recorded
\code{elapsed\_seconds} from the successful episode's process receipts.
These are supervised wall durations including capture, finalization,
and I/O; the timeout column is the declared cap, not an observation.
The outer run contains the paired native generation and four serial
graph judgments, together with modeled steps and orchestration. Its
28.36~s must not be added to those child durations. The separate saved-event
replay occurred later and is outside that outer run. No solver-only or
extraction-only time is inferred from the paired process.

\begin{table}[htbp]
\centering\small
\caption{Recorded process durations for one completed constructed
Lift T episode. Nested rows are components of the outer run;
these observations are not a latency or scaling benchmark.}
\label{tab:lift-episode-observed-times}
\begin{tabular}{@{}lrr@{}}
\toprule
Process & Observed wall (s) & Timeout (s)\\
\midrule
Outer modeled episode & 28.36 & 600\\
\quad Paired native generation & 10.82 & 30\\
\quad Released: production judge & 0.12 & 60\\
\quad Released: separate judge & 8.60 & 180\\
\quad Experimental: production judge & 0.15 & 60\\
\quad Experimental: separate judge & 8.35 & 180\\
Later saved-event replay & 0.11 & 120\\
\bottomrule
\end{tabular}
\end{table}

These measurements describe one completed run of the three-floor
example after the three retained technical UNKNOWN attempts below.
They are not a statistical sample, online adaptation latency,
hardware-normalized comparison, worst-case guarantee, or a scaling
result beyond the bounded models. Process-group RSS caps remain limits;
they are not measured aggregate-tree memory consumption.

The first observer attempt exceeded its 256-MiB stdout limit and
remains UNKNOWN; a separately recorded followup raised only stream
and output-tree limits to 512 and 768~MiB. The episode likewise retains
three earlier attempts as technical UNKNOWNs. Attempt~01 stopped when
the outer single-child supervisor detected a nested driver; the
followup called the existing driver mains in the episode process while
each native/checker child retained its own supervisor.
Attempt~02's stream capture hit the inherited outer 16-MiB file limit;
the next attempt raised that outer limit to 1~GiB.
Attempt~03 produced the complete event log and five valid inner
processes, but its final \code{wait4} RSS value of 557,793,280 bytes
exceeded the outer 512-MiB cap, so it remains resource UNKNOWN.
Attempt~04 raised the outer RSS cap to 2~GiB, retained all inner and
output limits, and completed; its independent episode audit passes.
Scientific inputs, candidate order, handoff, common input prefix and
acceptance criteria stayed fixed. Each followup has a separate plan,
output directory and receipt; no automatic retry rewrites an attempt.

The author-side records bind source, class, model, graph, checker,
process and event-log identities. Their entry points are
\path{CAUSAL_COMPARISON_CLOSURE.json},
\path{SUCCESS_CONTROLS_RESULT.json}, and
\path{adaptation_episode_04/EXECUTION_RECEIPT.json}, followed by
\path{adaptation_episode_04/independent/AUDIT_RESULT_attempt01.json}.
These distinguish the original released result, experimental result
and executed episode. The fourth public archive in Section~K
preserves them with explicit original/copy identities. Its consumer
replays saved semantics and the event audit; it does not regenerate
the native producer or execute a new adaptation episode.


\subsection*{Lift T: executed noninitial handoffs from a frozen cohort}
\label{supp:noninitial-handoff}
This separate followup executes existing native-controller prefixes to
noninitial source and monitor states before replacing the controller.
It preserves the preceding maintenance experiment and its distinct
fresh-generation protocol. The new experiment changes admission at all
ten declared boundaries; its twenty executed continuations all reach
fair lassos. The rejected candidates' cooperation defects are established
by complete rooted graphs, not by failures on these executed traces.


\paragraph{Construction, selection, and frozen denominator.}
The new contract is the same Lift T source used above. The constructed
old contract appends only the environment invariant $\G\neg b_3$;
its original initializers, system clauses, response monitors, and justice
lists are unchanged. Operationally, the old environment relation requires
both current and next $b_3$ to be false. The change removes exactly that
restriction at an atomic pause.

This design was chosen after inspecting the known Lift failure, saved
controllers, and all three single-button restrictions. Restricting $b_2$
still exposed a known non-one-hot native node. Restricting $b_1$ offered
unanswered-monitor states but lacked matching rank-zero targets for the
full proposed cohort and needed additional memory qualification.
That qualification and a separate genuinely pending cohort are supplied
in Section~R; the present historical experiment and its outcomes are unchanged.
The $b_3$ restriction supplied a small complete cohort and existing target
mappings. The restriction, every prefix, all four candidates, their order,
and both admission policies were then frozen before the new graph checks
and executions. This is an outcome-informed mechanism study on one known
source, not a held-out source or an independent benchmark.

From native node 0, the old mode reaches 26 private nodes. The cohort is
\emph{every} one of these nodes with a nonzero source response monitor:
ten private nodes representing eight source valuations. Equal source
values with different reached ranks remain separate rows. Both checking
routes independently verify the cohort's completeness. All ten source
states lie outside the original initializer and belong to the previously
certified original region $K$; membership is checked in the saved region,
without solving it again. These observations do not add contracts to the
eleven-replacement, fourteen-public-input, or 59-model populations, or
alter the separate eleven-state Elevator census and its 0/836 refinements.

\paragraph{Exact state and memory boundary.}
The ten retained source fields are $b_1,b_2,b_3,c$, $f_1,f_2,f_3$, and
\code{SYS\_CONSTRAINT.0.pRespondsToS.state} through
\code{SYS\_CONSTRAINT.2.pRespondsToS.state}. Write their values as
$(b,c,f,m)$, with each triple ordered by floor and $S0=0,S1=1$.
Every row has one $S1$ monitor whose corresponding floor response already
holds. The update
\[
 m_i'=(m_i\lor b_i)\land\neg(f_i\lor c)
\]
therefore clears that boundary monitor on the next transition; other
current requests can set other monitors. The experiment preserves
nonzero logical history at installation, but no unanswered obligation
already represented by an $S1$ monitor. Rows 03 and 07 do contain
unanswered current triggers whose corresponding monitors are still $S0$.

Table~\ref{tab:noninitial-roster} lists all source values and actual saved
private roots. The retained arm keeps the reached controller object,
node, and rank. A fresh arm creates a different private instance at a
matching node whose \emph{stored} rank is zero; it does not overwrite a
rank, reapply the source initializer, or reset a monitor. The source
object and all ten of its values remain unchanged during the switch.
Rows 09--10 retain rank zero already, but retention and fresh
initialization still denote different controller-instance operations.
Node IDs are local to each graph, even when their numbers agree.

\begin{table}[htbp]
\centering\small
\caption{Complete frozen noninitial roster. Old gives native node/rank;
N, Z, and P give actual rank-zero roots in the released, normalized,
and preference graphs. Every source valuation is noninitial.}
\label{tab:noninitial-roster}
\begin{tabular}{@{}ccrrrr@{}}
\toprule
ID & $(b,c,f,m)$ & Old & N & Z & P\\
\midrule
01 & $(010,0,010,010)$ & 728/1 & 626 & 94 & 626 \\
02 & $(010,1,010,010)$ & 729/1 & 627 & 95 & 627 \\
03 & $(110,0,010,010)$ & 730/1 & 646 & 102 & 646 \\
04 & $(110,1,010,010)$ & 731/1 & 647 & 103 & 647 \\
05 & $(100,0,100,100)$ & 656/2 & 492 & 68 & 492 \\
06 & $(100,1,100,100)$ & 657/2 & 493 & 69 & 493 \\
07 & $(110,0,100,100)$ & 716/2 & 580 & 84 & 580 \\
08 & $(110,1,100,100)$ & 717/2 & 581 & 85 & 581 \\
09 & $(010,0,010,010)$ & 626/0 & 626 & 94 & 626 \\
10 & $(010,1,010,010)$ & 627/0 & 627 & 95 & 627 \\
\bottomrule
\end{tabular}
\end{table}

\paragraph{Old-mode qualification and forty complete suffix checks.}
Before admission or runtime execution, both the source-equation route
and the saved-model route certify the old controller from both original
initial nodes, 0 and 1 (the two cancellation values). They agree on all
26 reached nodes and 120 edges and establish complete legal-input
coverage, system safety, strict liveness, and cooperation. A failure,
UNKNOWN, or disagreement at this gate would stop the direction without
changing the restriction or prefixes.

Each of the ten boundaries then receives four immutable graph/root
candidates in this fixed order: retain the old native instance; create
a fresh native instance; create a fresh normalized instance; create a
fresh preference instance. Each route checks all forty complete reachable
suffixes under the unchanged new source relation. These are explicitly
rooted judgments, not original-initializer judgments. At every reached
node, the check covers each original-legal next input with exactly one
saved safe response, preserving monitor updates and actual private
successors. Missing or ambiguous responses are UNKNOWN rather than an
invented edge. Strictness and cooperation use complete reachable-graph
justice checks, not bounded traces.

Both routes agree on all forty complete judgments in
Table~\ref{tab:noninitial-verdicts}; none is UNKNOWN. All candidates are
strict. The retained and fresh-native suffixes each have 154 nodes and
1,170 edges and fail cooperation, reaching native dead-end nodes 680 and
708. Each normalized and preference suffix has 154 nodes and 1,176 edges
and passes cooperation. For example, the native suffix from boundary 01
contains $728\to339\to438\to704\to680$; this graph witness was not the
executed continuation. Reinitializing native private memory alone does
not fix the rejected graph. The existing normalized alternative already
passes, so this experiment does not establish necessity of the preference.

\begin{table}[htbp]
\centering\small
\caption{All forty rooted candidate outcomes. Each cell gives
strict/cooperative verdicts (P: pass, F: fail). Both checking routes
agree, and all forty coverage and safety checks pass.}
\label{tab:noninitial-verdicts}
\begin{tabular}{@{}ccccc@{}}
\toprule
ID & Retain old & Fresh native & Fresh normalized & Fresh preference\\
\midrule
01 & P/F & P/F & P/P & P/P \\
02 & P/F & P/F & P/P & P/P \\
03 & P/F & P/F & P/P & P/P \\
04 & P/F & P/F & P/P & P/P \\
05 & P/F & P/F & P/P & P/P \\
06 & P/F & P/F & P/P & P/P \\
07 & P/F & P/F & P/P & P/P \\
08 & P/F & P/F & P/P & P/P \\
09 & P/F & P/F & P/P & P/P \\
10 & P/F & P/F & P/P & P/P \\
\bottomrule
\end{tabular}
\end{table}

\paragraph{Actual step, pause, admission, and continuation.}
For each boundary, each of the two policies creates an old native
instance at node 0 and actually calls its step/commit operations along
the fixed prefix in Table~\ref{tab:noninitial-executions}. It does not load
the endpoint as a substitute for execution. The same modeled plant and
source object then pause atomically, remove only $\G\neg b_3$, perform
admission, and resume. Strict-only admission accepts the first complete
strict candidate with coverage. Strict-and-cooperative admission also
requires cooperation. Both require the independently agreed exact
source/root/rank binding. If no candidate qualifies, the plant remains
paused; no such abstention occurs in these twenty executions.

At all ten boundaries, strict-only admission retains the reached native
instance. Requiring cooperation rejects both native candidates and
installs the existing normalized graph in a new private instance at the
specified rank-zero root. Thus all ten decisions differ and zero are
equal. The preference candidate qualifies in the certificates but is
never selected because normalization precedes it. There are forty
candidate judgments per checker and twenty runtime episodes; the fresh
native and preference candidates are not additional executed arms.

After admission, the fixed policy offers the lexicographically smallest
source-legal next input in $(b_1,b_2,b_3,c)$ order, false before true.
It enumerates the sixteen typed input valuations and checks legality
before committing a unique saved response. Each suffix is bounded by
32 transitions and stops at the first repeated complete private state,
dead end, rejected step, or UNKNOWN. A repeated state is a fair lasso
only when its cycle meets every justice: Lift has no original environment
justices and requires each response monitor to be $S0$ infinitely often.

\begin{table}[htbp]
\centering\small
\caption{All twenty executed outcomes. Each row executes the old prefix
twice, once per policy. The second and third sequences begin at the
installed root and end with an actually traversed self-loop. Every
suffix is a fair lasso; ID numbers refer to Table~\ref{tab:noninitial-roster}.}
\label{tab:noninitial-executions}
\begin{tabular}{@{}cL{0.31\linewidth}L{0.25\linewidth}L{0.25\linewidth}@{}}
\toprule
ID & Old native prefix & Strict-only suffix & Cooperative suffix\\
\midrule
01 & $0\to8\to728$ & $728\to336\to336$ & $94\to16\to16$ \\
02 & $0\to8\to729$ & $729\to336\to336$ & $95\to16\to16$ \\
03 & $0\to8\to730$ & $730\to656\to0\to0$ & $102\to104\to0\to0$ \\
04 & $0\to8\to731$ & $731\to336\to336$ & $103\to16\to16$ \\
05 & $0\to8\to730\to656$ & $656\to0\to0$ & $68\to0\to0$ \\
06 & $0\to8\to730\to657$ & $657\to0\to0$ & $69\to0\to0$ \\
07 & $0\to10\to410\to716$ & $716\to626\to336\to336$ & $84\to94\to16\to16$ \\
08 & $0\to10\to410\to717$ & $717\to336\to336$ & $85\to16\to16$ \\
09 & $0\to10\to410\to716\to626$ & $626\to336\to336$ & $94\to16\to16$ \\
10 & $0\to10\to410\to716\to627$ & $627\to336\to336$ & $95\to16\to16$ \\
\bottomrule
\end{tabular}
\end{table}

All twenty executions complete: 56 old-prefix commits and 44 post-change
commits, recorded in 302 events. Each suffix reaches a fair self-loop
in two steps, except both policies at boundaries 03 and 07, which take
three. There are no executed failures, incomplete outcomes, or UNKNOWNs.
No post-change offered input sets $b_3=1$. Consequently these traces show
successful bounded execution and different admission decisions, but no
observed failure avoidance, improvement in service success, or activation
of a newly allowed third-floor request. A retained graph can have the
observed fair lasso while still failing cooperation elsewhere. The
complete rooted certificates establish that difference.

\paragraph{Independent checks, qualification, and trust scope.}
The source route derives legal inputs, system safety, and monitor updates
from separately written source equations, using its own graph analysis.
The saved-model route decodes the pinned BDD exports and independently
checks typed source equality before its graph analysis. Both retain the
exported identities, original-initial markers, and actual ranks; native
printer parity is checked where a printer exists. The normalized export
has no native printer. Both routes share the exported candidate graphs
and provenance, and their code was mutually visible; this is not a blind
implementation experiment or proof of compiler independence.

A separate source-based event auditor imports neither the runtime nor
its transition implementation. It reconstructs every old prefix,
step/commit, pause, admission, installation, source update, private
successor, and fair cycle. All twenty episodes pass, including exact
source-object continuity and preserved nonzero boundary monitors. Outer
checks bind the event PID to the owned child, check event sequence and
clock order, and reconcile the process receipts. All six scientific
processes complete once and are reaped with their process groups absent:
two old-mode checks, two forty-root checks, one twenty-episode execution,
and one event audit. No new synthesis, native extraction, compiler run,
or region solution occurs.

Before each relevant stage, synthetic qualification includes 30 source-checker cases,
55 saved-model-checker cases, 37 runtime cases, 21 event-auditor unit
cases, and 12 full-envelope cases, plus full twenty-episode positive and
abstention fixtures. The latter establish the audit's handling of both
installation and remaining paused, not additional scientific episodes.
Code revisions and failed synthetic attempts remain recorded. The
scientific cohort is unchanged and has one attempt per stage; no failed
scientific row is omitted or retried.

The frozen bounds are 180 seconds per checker and 120 seconds for the
runtime, at most two concurrent children, and a sampled 2-GiB RSS limit
per child. The memory threshold is not an operating-system address-space
cap. Timeouts or exceeded limits produce UNKNOWN with cleanup of the
owned process group. These bounds are not measured adaptation latency or
scaling results; the preceding maintenance episode's 28.36-second
measurement does not describe this saved-graph experiment.

\paragraph{Exact input identity and saved evidence.}
Table~\ref{tab:noninitial-input-hashes} records source and target graph
identities of the frozen author-side inputs; the public manifest records
original/copy identities when paths are rewritten. The plan and
stage-specific execution authorizations bind model predicates, printer
exports, the existing $K$ result, code, and declared bounds. Retention
and fresh-native initialization use the same graph bytes with distinct
root/instance policies. The constructed old source is the unchanged
new source followed only by the declared $\G\neg b_3$ assumption.

{\small
\begin{longtable}{@{}L{0.23\linewidth}rL{0.62\linewidth}@{}}
\caption{Exact frozen input identities (full SHA-256).}
\label{tab:noninitial-input-hashes}\\
\toprule
Input & Bytes & SHA-256\\
\midrule
\endfirsthead
\toprule
Input & Bytes & SHA-256\\
\midrule
\endhead
New Lift source & 2218 & \path{6a8f8220385130b2eb219ecfa33186b179d610ffcf805a30f9cab47ab204dae7} \\
Constructed old source & 2277 & \path{8b32937fd7b465d1a1052ca5a2bb576cac686383a7d54015f44fa21442ff88f5} \\
Released graph & 860323 & \path{bbee6ff9a32c55691f9b7b58c162e4e2c2af02d22792967537751ccb9d148a1f} \\
Normalized graph & 97767 & \path{50ea0e539eea65b0fec615200fd0dffa973caabf90d515421aa358556caa3ca4} \\
Preference graph & 843906 & \path{b478d31d20641af7c0e97a17e9d7225d95dadb3880f2562bf889a7405d8fbba7} \\
Frozen plan & 35107 & \path{43ff75e0562bb6c554eb018917c78cd1cc721f1b5e18db43582cbbb3a8b22539} \\
Frozen roster & 12579 & \path{815e9484e48ae29813deb5077a6272da9ae9f77215262c0b19a4168d65614b5e} \\
\bottomrule
\end{longtable}
}

The eighth public static snapshot~\cite{savedartifact} is identified by
\path{noninitial-handoff-evidence-20261021-v1}, with
\path{NONINITIAL_HANDOFF.md} and \path{NONINITIAL_HANDOFF_MANIFEST.json}.
Its scientific entry records are \path{evidence/EXPERIMENT_PLAN.json},
\path{evidence/COHORT.json}, \path{evidence/ROOT_OLD_GATE_ADOPTION.json},
\path{evidence/ROOT_CANDIDATE_CERTIFICATES_ADOPTION.json},
\path{execution/runtime_01/RESULT.json},
\path{execution/runtime_01/EVENTS.jsonl}, and
\path{execution/INDEPENDENT_EVENT_AUDIT.json}. The derived
\path{SUMMARY.json} records the scientific adoption's observations and
limits with its original hash; the private coordination record itself
is omitted. These distinguish actual author-side execution from
inspection of saved evidence. The older public
regeneration/replay entry point does not execute or replay these new
handoffs; its original counts retain their original scope.

This followup narrows evidence about executed noninitial source and
monitor continuity with explicit target-memory initialization. It does
not establish arbitrary state or private-memory migration, carryover of
unanswered obligations already represented by $S1$ monitors, unseen
failure coverage, asynchronous or hardware
deployment, or an additional active extraction rescue. Its admission
difference uses rooted cooperation certificates; it supplies no new
practical decision requiring the exact $W_C/W_F$ strict-preservation
conditions and no new refusal transfer. The existing Vehicle and
preservation evidence remain separate.


\paragraph{Lift P: the whole-model timing mismatch.}
Let $B=b_1\land b_2\land b_3$, $\phi=\neg c\lor\neg B$, and
$\psi=f_1\lor\neg B$. The added source assumptions give the documented
original translation
\[
\rho_e^P=E_L\land\phi(b',c')\land\psi(b,f),\qquad
\theta_e^P=(\neg b_1\land\neg b_2\land\neg b_3)\land\phi(b,c).
\]
The mixed invariant $\psi$ is evaluated at the current state; the
environment-only $\phi$ is evaluated initially and at the next input.
Both initial cancellation values remain allowed. The reduced helper-false
relation instead checks $\phi$ at the current input. The remaining environment clauses agree with the
independently encoded source equations. With both helper values fixed
false, all 2,048 assignments to the eleven original bits in the union
of environment-relation supports are checked. The relations differ
on 32: 24 false-to-true and eight true-to-false.

One witness starts at the original joint initial state
$(0,0,0,0;1,0,0;S0,S0,S0)$ and moves to
$(1,1,1,1;1,0,0;S0,S0,S0)$ with helper false at both ends.
Original environment safety is false, reduced environment safety is
true, and both system relations are true. A control with the opposite
environment-relation disagreement is also checked. Initializers and ordered
justices pass their separate checks. This diagnosis exhausts the
environment relation's support, not the full system-relation domain;
system safety is checked at the declared witnesses. This blocks transfer
of the native verdict at the whole-model correspondence gate. The earlier
diagnostic export-schema failure is retained as an integration failure,
not a semantic witness. A separate selected-graph audit follows.

\paragraph{Post-hoc direct original-contract graph certificate.}
We separately audit the same saved native graph using the original
predicates, without transferring its native verdict. The view begins at
both original initial nodes and retains helper-false responses to every
original-legal input, with raw node identities, private ranks, and AUX
values unchanged. All 16 input valuations are checked at each of its
136 reachable states. The 958 legal-input responses are unique and satisfy
original SYS/AUX safety; no missing response is repaired and no unsafe
legal edge is discarded. The resulting 958-edge view passes strictness
but fails cooperation. All 16 failing states are environmental dead ends
with zero original-legal inputs, rather than missing controller responses.

Hand-derived source predicates with bitset transitive closure agree with
a separate original-BDD, value-decoding, and Kosaraju/reverse-reachability
route on all saved states, edges, inputs, ordered justices, and SCC evidence.
Both routes share the frozen graph and prior source qualification.
The BDD route reuses the earlier saved-BDD decoder, and implementation
code was visible across routes; this is not independent compilation
or a blind implementation comparison.
This certificate leaves the whole-model mismatch and the historical
correspondence-gated table unchanged. It adds no native synthesis run,
new candidate, or initially cooperative false-refusal example.

\paragraph{A source obstruction independent of the selected graph.}
From either legal initial input, $b=000$, $f=100$, and all monitors are
$S0$. The environment chooses $(b',c')=(010,0)$.
Source system safety forces $f_1'=1$. Any additional true floor immediately
violates the environment's current one-floor assumption; otherwise the
next state has $b=010$, $c=0$, $f=100$, and all monitors $S0$.
The environment can then choose $(b',c')=(111,0)$: button persistence,
button clearing, and the next-input invariant all permit this move.
System safety forces $f_2'=1$. An additional true floor violates the
one-floor assumption; with only $f_2'$ true, the next state instead
violates the mixed current-state assumption
$f_1\lor\neg(b_1\land b_2\land b_3)$.
Thus every safe response yields an environmental dead end within two
rounds. A strict controller that avoids the first dead end must answer
the second legal input safely; hence no strict cooperative strategy
exists from either original initial input,
regardless of private memory. This written source argument is independent
of the selected graph and does not recompute a winning region; it does
not claim that the intermediate state lacks an existential continuation.
The argument uses the documented translation of mixed ENV/SYS invariants
as current-state environment safety. Initially $b=000$, and both chosen
inputs have $c'=0$, so $\phi$ is true at every current and next valuation
used here; its current/next timing does not affect this obstruction.

\section*{I. A declared pair separating region membership from selection}
\label{supp:choice-pair}
The pair was fixed before native execution and is separate from the
eleven published candidates. Choice has one environment Boolean $e$,
one system Boolean $p$, no logical auxiliary, and
\[
\begin{aligned}
\theta_e&=\neg e,&\theta_s&=\neg p,&\rho_e&=(e'\Rightarrow p),\\
\rho_s&=\top,& A&=\mathbf{G}\mathbf{F}e,&G&=\mathbf{G}\mathbf{F}\neg p.
\end{aligned}
\]
Toggle changes only system safety to $p'\leftrightarrow\neg p$.

\paragraph{Original source proof.}
From every one of the four original typed states, policy $p'=\neg p$
satisfies system safety and $G$ on every play in either contract.
After every finite original-legal prefix, the environment can choose
$e'=p$ forever; this is legal and makes $e$ alternate. Thus this same
policy is strict and cooperative from every state, yielding
$W=V=K=C=D$ and nonvacuous $I_M(K)=\mathit{true}$ for both contracts.
In Choice, policy $p'=0$ instead satisfies $G$ itself but forces $e'=0$
after reaching $p=0$, so no assumption-fair suffix remains. Every
state of that bad policy still belongs to $K$.

\paragraph{Audited complete transformed relations.}
For helper $h$, both saved native models satisfy
\[
\begin{aligned}
\theta_e^R&=\neg e\land\neg h,&\theta_s^R&=\neg p,\\
\rho_e^R&=((e'\Rightarrow p)\lor h')\land(\neg h\lor h'),&
\rho_s^R&=\rho_s\lor h,\\
A^R&=\mathbf{G}\mathbf{F}e\land\mathbf{G}\mathbf{F}\neg h,&G^R&=G.
\end{aligned}
\]
An independent saved-DAG check establishes these equations on all eight
typed native states and all 64 current/next pairs in each case.
Field ownership, domains and initialization are checked on the full
typed native model; justice identities and order are checked at the
original, preconstructor and solver phases. Hence
the same $p$-toggle policy satisfies native system safety and $G$
unconditionally, and its helper-false restriction has the cooperative
source behavior just proved. The transformation does not exclude this
good original-view policy. This does not demand cooperation in the full
transformed arena after an environment-chosen helper flip.

\paragraph{Actual selected outputs.}
Choice returns three native states and seven edges. Its original view
has only node 0, with $(e,p)=(0,0)$, rank 0, and a self-loop. It satisfies
$G$ and fails cooperation while remaining inside $K=V=D$.
Toggle returns seven native states and 18 edges. Its original view has
nodes $0=(0,0)$, $1=(0,1)$ and $6=(1,0)$, all rank 0, with edges
$0\to1$, $1\to0$, $1\to6$ and $6\to1$. It passes both checks.
Each case received one native attempt; complete typed source
correspondence and independent graph checks succeeded for both. Saved
state, edge, justice, SCC and raw-node comparisons agree exactly.

The pair shows that $I_M(K)$ and membership of every reached state in
$K$ do not certify the particular output's cooperation. It does not
attribute that selection separately to the raw fixed-point core or
the concrete constructor. Both unnormalized original arenas admit
legal $G$-fair, $A$-unfair plays, so they do not refute the conditional
language-inclusion theorem of~\cite{friendly}. The pair adds no public
candidate, learned repair, adaptation episode or deployment result.

